\section{Glossário}
\begin{itemize}
    \item \textbf{Agendamento:} Ação de reservar um bloco de tempo (Slot) na agenda de um profissional de saúde para o atendimento de um paciente específico.
    \item \textbf{Atendimento:} O evento no qual a interação direta entre o Paciente e o Profissional de Saúde de fato ocorre, sendo o momento onde as observações e diagnósticos básicos são registrados no sistema.
    \item \textbf{CRM (Customer Relationship Management):} Refere-se à plataforma voltada para o gerenciamento do relacionamento com o paciente, focando no histórico de contatos, otimização da agenda, comunicação de lembretes e retenção.
    \item \textbf{FHIR (Fast Healthcare Interoperability Resources):} Padrão desenvolvido pelo HL7 para troca eletrônica de informações de saúde.
    \item \textbf{LGPD (Lei Geral de Proteção de Dados):} Legislação brasileira que regula o tratamento de dados pessoais.
    \item \textbf{Paciente:} Indivíduo cadastrado na clínica que recebe ou está agendado para receber cuidados médicos ou administrativos. É a entidade central em torno da qual o CRM opera.
    \item \textbf{Profissional de Saúde:} Colaborador da clínica autorizado a realizar os atendimentos clínicos, visualizar as agendas de consultas e registrar anotações no histórico de saúde do paciente.
\end{itemize}